\begin{figure}[ht]
			\centering
			\resizebox{\textwidth}{!}{%
				\begin{tikzpicture}[>=Stealth, line width=0.8pt]
					
					\def\R{1.4}
					\def\arr{0.55}
					\def\ss{0.65}
					\def\xV{-2.8}      % retta r_V a sinistra dell'arco
					\def\yH{-1.3}
					\def\rt{0.28}
					\def\dxp{7.0}      % separazione tra i due pannelli (da rifinire)
					\def\rA{1.0}
					
					
					% ============================================================
					% PANNELLO (a) — Cedimento s_1 (orizzontale in A)
					% ============================================================
					\begin{scope}[xshift=0cm]
						
						% --- semiarchi ---
						\draw[line width=1.2pt] (-\R,0)
						arc[start angle=180, end angle=90, radius=\R];
						\draw[line width=1.2pt] (0,\R)
						arc[start angle=90, end angle=0, radius=\R];
						
						% --- cerniera al suolo in A ---
						\draw[thick] (-\R,0) -- (-\R-0.18,-0.30)
						-- (-\R+0.18,-0.30) -- cycle;
						\draw[thick] (-\R-0.28,-0.30) -- (-\R+0.28,-0.30);
						\foreach \x in {-0.20,-0.10,0,0.10,0.20}{
							\draw[thin] (-\R+\x,-0.30) -- (-\R+\x-0.09,-0.42);
						}
						% --- cerniera al suolo in B ---
						\draw[thick] (\R,0) -- (\R-0.18,-0.30)
						-- (\R+0.18,-0.30) -- cycle;
						\draw[thick] (\R-0.28,-0.30) -- (\R+0.28,-0.30);
						\foreach \x in {-0.20,-0.10,0,0.10,0.20}{
							\draw[thin] (\R+\x,-0.30) -- (\R+\x-0.09,-0.42);
						}
						
						% --- perni in A, B (esterni) e D (interno) ---
						\node[circle, fill=white, draw=black, inner sep=0pt,
						minimum size=4pt, line width=0.7pt] at (-\R,0) {};
						\node[circle, fill=white, draw=black, inner sep=0pt,
						minimum size=4pt, line width=0.7pt] at (\R,0) {};
						\node[circle, fill=white, draw=black, inner sep=0pt,
						minimum size=4pt, line width=0.7pt] at (0,\R) {};
						
						% --- costruzione di C_1: verticale per A e retta BD ---
						\draw[gray!55, dashed, thin] (-\R,-0.50) -- (-\R, 2*\R+0.40);
						\draw[gray!55, dashed, thin]
						({\R+0.25},-0.25) -- ({-\R-0.20}, 2*\R+0.20);
						
						% --- centro C_1 (label spostata a destra: la sinistra è
						\fill (-\R, 2*\R) circle (1.8pt);
						\node[font=\scriptsize, anchor=west] at ({-\R+0.10}, 2*\R) {$C_1$};
						
						% --- etichette punti notevoli ---
						\node[font=\scriptsize, anchor=east] at ({-\R-0.10}, 0.18) {$A$};
						\node[font=\scriptsize, anchor=west] at ({\R}, 0.) {$B\!\equiv\!C_2$};
						\node[font=\scriptsize, anchor=south] at (0.0, \R){$D\!\equiv\!C_{12}$};
						
						% --- etichette dei corpi (interno della rispettiva cupola) ---
						\node[font=\scriptsize]	at ({0.95*cos(110)-0.3},{0.95*sin(110)}) {$\mathcal{C}_1$};
						\node[font=\scriptsize]	at ({0.95*cos(70)+1.0},{0.95*sin(70)})	{$\mathcal{C}_2$};
						
						% --- cedimento s_1 in A ---
						\draw[->, thick] (-\R,0) -- ({-\R+\ss},0);
						\node[font=\scriptsize, anchor=south] at ({-\R+\ss/2}, 0.04) {$s_1$};
						
						%			
						% ==========================================================
						% retta r_V (verticale, a sinistra dell'arco) — campo u(y)
						% ==========================================================
						\draw[gray!45, thin] (\xV, -0.30) -- (\xV, 2*\R+0.40)
						node[above, font=\scriptsize] {$r_V$};
						
						% proiezioni dei punti notevoli su r_V
						\draw[gray!40, dashed, thin] (-\R, 0) -- (\xV, 0);
						\draw[gray!40, dashed, thin] (0, \R) -- (\xV, \R);
						\draw[gray!40, dashed, thin] (-\R, 2*\R) -- (\xV, 2*\R);
						
						\fill[gray!75] (\xV, 0) circle (1.3pt);
						\fill[gray!75] (\xV, \R) circle (1.3pt);
						\fill[gray!75] (\xV, 2*\R) circle (1.3pt);
						
						% --- skyline u_1^(1) di C_1 (gray solido) ---
						\draw[gray!75, thin]
						({\xV+\ss}, 0) -- (\xV, 2*\R);
						\foreach \k in {0,1,2,3,4}{
							\pgfmathsetmacro{\yy}{\k*\R/4}
							\pgfmathsetmacro{\uu}{\ss*(1-\yy/(2*\R))}
							\draw[->, gray!80, thin]
							(\xV,\yy) -- ({\xV+\uu},\yy);
						}
						\node[gray!85, font=\scriptsize, anchor=north west]
						at ({\xV+\ss-0.2}, 1.0) {$u_1$};
						
						% --- skyline u_2 di C_2 (nero tratteggiato per distinguerlo) ---
						\draw[black!55, thin, dashed]
						(\xV, 0) -- ({\xV+\ss/2}, \R);
						\foreach \k in {1,2,3,4}{
							\pgfmathsetmacro{\yy}{\k*\R/4}
							\pgfmathsetmacro{\uu}{(\ss/2)*\k/4}
							\draw[->, black!60, thin, dashed]
							(\xV,\yy) -- ({\xV+\uu},\yy);
						}
						\node[black!70, font=\scriptsize, anchor=east] at ({\xV}, \R-0.7) {$u_2$};
						
						
						% --- etichette delle proiezioni dei centri su r_V ---
						\node[font=\scriptsize, anchor=north east]	at ({\xV-0.05}, {2*\R+0.2}) {$C_1$};
						\node[font=\scriptsize, anchor=north east]	at ({\xV-0.05}, {\R+0.3}) {$C_{12}$};
						\node[font=\scriptsize, anchor=north east]	at ({\xV-0.05}, 0.2) {$C_2$};
						
						
						
						% ==========================================================
						% retta r_H (orizzontale) — campo v(x)
						% ==========================================================
						\draw[gray!45, thin] (-\R-0.40, \yH) -- (\R+0.40, \yH)
						node[right, font=\scriptsize] {$r_H$};
						
						\draw[gray!40, dashed, thin] (-\R, -0.45) -- (-\R, \yH);
						\draw[gray!40, dashed, thin] (0, \R-0.10) -- (0, \yH);
						\draw[gray!40, dashed, thin] (\R, -0.45) -- (\R, \yH);
						
						\fill[gray!75] (-\R, \yH) circle (1.3pt);
						\fill[gray!75] (0, \yH) circle (1.3pt);
						\fill[gray!75] (\R, \yH) circle (1.3pt);
						
						% --- skyline v_1 di C_1 ---
						\draw[gray!75, thin]
						(-\R, \yH) -- (0, {\yH+\ss/2});
						\foreach \k in {1,2,3,4}{
							\pgfmathsetmacro{\xx}{-\R+\k*\R/4}
							\pgfmathsetmacro{\vv}{(\ss/2)*\k/4}
							\draw[->, gray!80, thin]
							(\xx,\yH) -- (\xx, {\yH+\vv});
						}
						\node[gray!85, font=\scriptsize] at ({-\R/2-0.05}, {\yH-0.30}) {$v_1$};
						
						\node[gray!85, font=\scriptsize, anchor=south] at ({\xV+\ss/2}, -0.4) {$s_1$};
						
						% --- skyline v_2 di C_2 ---
						\draw[gray!75, thin]
						(0, {\yH+\ss/2}) -- (\R, \yH);
						\foreach \k in {0,1,2,3}{
							\pgfmathsetmacro{\xx}{\k*\R/4}
							\pgfmathsetmacro{\vv}{(\ss/2)*(1-\k/4)}
							\ifnum\k>0
							\draw[->, gray!80, thin]
							(\xx,\yH) -- (\xx, {\yH+\vv});
							\fi
						}
						\node[gray!85, font=\scriptsize]
						at ({\R/2+0.05}, {\yH-0.30}) {$v_2$};
						
						\node[font=\scriptsize, anchor=south] at (0.0, {\yH+\ss/3}) {$s_1/2$};
						
						
						
						% --- etichette delle proiezioni dei centri su r_H ---
						\node[font=\scriptsize, anchor=north] at (-\R, {\yH-0.04}) {$C_1$};
						\node[font=\scriptsize, anchor=north]	at (0, {\yH-0.04}) {$C_{12}$};
						\node[font=\scriptsize, anchor=north] at (\R, {\yH-0.04}) {$C_2$};
						
						
						% Configurazione variata (qualitativa, schematica)
						\draw[gray!65, thin] (-\R, 2*\R) -- ({-\R+\ss}, 0);
						\draw[gray!65, thin] (-\R, 2*\R) -- ({\ss/2}, {\R+\ss/2});
						\draw[gray!65, thin] (\R, 0) -- ({\ss/2}, {\R+\ss/2});
						
						\fill[gray!65] ({\ss/2}, {\R+\ss/2}) circle (1.5pt);
						
						\node[gray!65,font=\scriptsize, anchor=west] at ({-\R+\ss-0.08}, 0.0) {$A'$};
						\node[gray!65,font=\scriptsize, anchor=south west]
						at ({\ss/2+0.05}, {\R+\ss/2+0.05}) {$D'$};
						
						% --- semiarchi deformati (linearizzazione del campo di rotazione) ---
						\draw[gray!65, line width=1pt]
						plot[domain=90:180, samples=30, smooth]
						({\R*cos(\x) + \ss*(1 - sin(\x)/2)},
						{\R*sin(\x) + (\ss/2)*(cos(\x) + 1)});
						\draw[gray!65, line width=1pt]
						plot[domain=0:90, samples=30, smooth]
						({\R*cos(\x) + (\ss/2)*sin(\x)},
						{\R*sin(\x) + (\ss/2)*(1 - cos(\x))});
						
						% --- archi degli angoli di rotazione ---
						\pgfmathsetmacro{\dtheta}{atan(\ss/(2*\R))}
						\def\rA{1.0}
						% vartheta_1 fra verticale per A e C_1A'  (antiorario)
						\draw[->, gray!75, thin]
						({-\R + \rA*cos(270)}, {2*\R + \rA*sin(270)})
						arc[start angle=270, end angle={270+\dtheta}, radius=\rA];
						% vartheta_1 fra C_1D (= retta BD prolungata) e C_1D'  (antiorario)
						\draw[->, gray!75, thin]
						({-\R + \rA*cos(315)}, {2*\R + \rA*sin(315)})
						arc[start angle=315, end angle={315+\dtheta}, radius=\rA];
						% vartheta_2 fra BD e BD'  (orario)
						\draw[->, gray!75, thin]
						({\R + \rA*cos(135)}, {\rA*sin(135)})
						arc[start angle=135, end angle={135-\dtheta}, radius=\rA];
						
						% --- etichette degli angoli, accanto alla punta della freccia ---
						\node[gray!85, font=\scriptsize, anchor=west]
						at ({-\R + \rA*cos(270+\dtheta) - 0.08}, {2*\R + \rA*sin(270+\dtheta)})
						{$\vartheta_1$};
						\node[gray!85, font=\scriptsize, anchor=west]
						at ({-\R + \rA*cos(315+\dtheta) - 0.08}, {2*\R + \rA*sin(315+\dtheta)+0.08})
						{$\vartheta_1$};
						\node[gray!85, font=\scriptsize, anchor=east]
						at ({\R + \rA*cos(135-\dtheta) - 0.06}, {\rA*sin(135-\dtheta)})
						{$\vartheta_2$};
						
						% --- etichetta pannello (centrata sull'estensione effettiva) ---
						\node[font=\scriptsize] at ({(\xV-\R)/2}, \yH-0.55) {(a)};
						
					\end{scope}
					
					% ============================================================
					% PANNELLO (b) — Cedimento s_2 (verticale in A)
					% ============================================================
					\begin{scope}[xshift=\dxp cm]   % posizione del pannello (b)
						
						% --- semiarchi ---
						\draw[line width=1.2pt] (-\R,0)
						arc[start angle=180, end angle=90, radius=\R];
						\draw[line width=1.2pt] (0,\R)
						arc[start angle=90, end angle=0, radius=\R];
						
						% --- cerniera al suolo in A ---
						\draw[thick] (-\R,0) -- (-\R-0.18,-0.30)
						-- (-\R+0.18,-0.30) -- cycle;
						\draw[thick] (-\R-0.28,-0.30) -- (-\R+0.28,-0.30);
						\foreach \x in {-0.20,-0.10,0,0.10,0.20}{
							\draw[thin] (-\R+\x,-0.30) -- (-\R+\x-0.09,-0.42);
						}
						% --- cerniera al suolo in B ---
						\draw[thick] (\R,0) -- (\R-0.18,-0.30)
						-- (\R+0.18,-0.30) -- cycle;
						\draw[thick] (\R-0.28,-0.30) -- (\R+0.28,-0.30);
						\foreach \x in {-0.20,-0.10,0,0.10,0.20}{
							\draw[thin] (\R+\x,-0.30) -- (\R+\x-0.09,-0.42);
						}
						
						% --- perni in A, B (esterni) e D (interno) ---
						\node[circle, fill=white, draw=black, inner sep=0pt,
						minimum size=4pt, line width=0.7pt] at (-\R,0) {};
						\node[circle, fill=white, draw=black, inner sep=0pt,
						minimum size=4pt, line width=0.7pt] at (\R,0) {};
						\node[circle, fill=white, draw=black, inner sep=0pt,
						minimum size=4pt, line width=0.7pt] at (0,\R) {};
						
						
						% --- costruzione di C_1: orizzontale per A e retta BD ---
						\draw[gray!55, dashed, thin] (-\R-0.30, 0) -- (\R+0.30, 0);
						\draw[gray!55, dashed, thin] (\R+0.25, -0.25) -- (-0.40, \R+0.40);
						
						% --- etichette punti notevoli ---
						\node[font=\scriptsize, anchor=east] at ({-\R-0.10}, 0.18) {$A$};
						\node[font=\scriptsize, anchor=west] at ({\R+0.10}, 0.18){$B\!\equiv\!C_1\!\equiv\!C_2$};
						\node[font=\scriptsize, anchor=south] at (0.0, \R){$D\!\equiv\!C_{12}$};
						
						% --- etichette dei corpi (interno della rispettiva cupola) ---
						\node[font=\scriptsize]	at ({0.95*cos(110)-1.2},{0.95*sin(110)}) {$\mathcal{C}_1$};
						\node[font=\scriptsize]	at ({0.95*cos(70)+1.20},{0.95*sin(70)})	{$\mathcal{C}_2$};
						
						% --- cedimento s_2 in A (verticale, verso il basso) ---
						\draw[->, thick] (-\R, 0) -- (-\R, -\ss);
						\node[font=\scriptsize, anchor=east] at ({-\R-0.32}, -\ss/2) {$s_2$};
						
						% --- costruzione di C_1: orizzontale per A e retta BD ---
						\draw[gray!55, dashed, thin] (-\R-0.30, 0) -- (\R+0.30, 0);
						\draw[gray!55, dashed, thin] (\R+0.25, -0.25) -- (-0.40, \R+0.40);
						
						% ============================================================
						% retta r_V (verticale, a sinistra) — campo u(y)
						% ============================================================
						\draw[gray!45, thin] (\xV, -0.30) -- (\xV, \R+0.40)	node[above, font=\scriptsize] {$r_V$};
						
						% proiezioni dei punti notevoli su r_V
						\draw[gray!40, dashed, thin] (\R, 0) -- (\xV, 0);
						\draw[gray!40, dashed, thin] (0, \R) -- (\xV, \R);
						
						\fill[gray!75] (\xV, 0) circle (1.3pt);
						\fill[gray!75] (\xV, \R) circle (1.3pt);
						
						% etichette delle tracce — a destra di r_V (skyline a sx)
						\node[font=\scriptsize, anchor=east]	at ({\xV}, 0.0) {$C_1\!\equiv\!C_2$};
						\node[font=\scriptsize, anchor=west]	at ({\xV+0.05}, {\R-0.04}) {$C_{12}$};
						
						% skyline u(y) — L-shape a sinistra di r_V
						\draw[gray!75, thin]
						(\xV, 0) -- ({\xV-\ss/2}, \R) -- (\xV, \R);
						\foreach \k in {1,2,3,4}{
							\pgfmathsetmacro{\yy}{\k*\R/4}
							\pgfmathsetmacro{\uu}{\k*\ss/8}
							\draw[->, gray!75, thin] (\xV, \yy) -- ({\xV-\uu}, \yy);
						}
						
						% picco label "s_2/2" al midpoint freccia y=R
						\node[font=\scriptsize, anchor=south] at ({\xV-\ss/4}, {\R}) {$s_2/2$};
						
						% nome del campo, sotto la base del r_V
						\node[font=\scriptsize, anchor=north east] at ({\xV-\ss/4}, {\R/2+0.4}) 	{$u_1\!=\!u_2$};
						
						% ============================================================
						% retta r_H (orizzontale, sotto) — campo v(x)
						% ============================================================
						\draw[gray!45, thin] (-\R-0.30, \yH) -- (\R+0.40, \yH)
						node[right, font=\scriptsize] {$r_H$};
						
						% proiezioni dei punti notevoli su r_H
						\draw[gray!40, dashed, thin] (-\R, 0) -- (-\R, \yH);
						\draw[gray!40, dashed, thin] (0, \R) -- (0, \yH);
						\draw[gray!40, dashed, thin] (\R, 0) -- (\R, \yH);
						
						\fill[gray!75] (-\R, \yH) circle (1.3pt);
						\fill[gray!75] (0, \yH) circle (1.3pt);
						\fill[gray!75] (\R, \yH) circle (1.3pt);
						
						% etichette delle tracce — sopra r_H (skyline sotto)
						\node[font=\scriptsize, anchor=south]	at (-\R, {\yH+0.04}) {$A$};
						\node[font=\scriptsize, anchor=south]	at (0, {\yH+0.04}) {$C_{12}$};
						\node[font=\scriptsize, anchor=south]	at (\R, {\yH+0.04}) {$C_1\!\equiv\!C_2$};
						
						% skyline v(x) — L-shape sotto r_H
						\draw[gray!75, thin]
						(-\R, \yH) -- (-\R, {\yH-\ss}) -- (\R, \yH);
						\foreach \k in {0,1,2,3,4,5,6,7}{
							\pgfmathsetmacro{\xx}{-\R + \k*\R/4}
							\pgfmathsetmacro{\vv}{\ss*(1-\k/8)}
							\draw[->, gray!75, thin] (\xx, \yH) -- (\xx, {\yH-\vv});
						}
						
						% picco label "s_2" al midpoint freccia x=-R
						\node[font=\scriptsize, anchor=east] at ({-\R+0.04}, {\yH-\ss/2}) {$s_2$};
						
						% nome del campo, sotto e a sx della picco
						\node[font=\scriptsize, anchor=north east]	at ({\R/2-0.05}, {\yH-\ss+0.2}) {$v_1\!=\!v_2$};
						
						% ============================================================
						% configurazione variata: rotazione rigida attorno a B
						% ============================================================
						
						% raggi della configurazione variata
						\draw[gray!65, thin] (\R, 0) -- (-\R, -\ss);             % BA'
						\draw[gray!65, thin] (\R, 0) -- ({-\ss/2}, {\R-\ss/2});  % BD'
						
						% semiarchi deformati (linearizzazione del campo di rotazione)
						\draw[gray!75, line width=0.9pt]
						plot[domain=90:180, samples=24, smooth, variable=\t]
						({\R*cos(\t) - \ss/2*sin(\t)}, {\R*sin(\t) + \ss/2*(cos(\t)-1)});
						\draw[gray!75, line width=0.9pt]
						plot[domain=0:90, samples=24, smooth, variable=\t]
						({\R*cos(\t) - \ss/2*sin(\t)}, {\R*sin(\t) + \ss/2*(cos(\t)-1)});
						
						% archi degli angoli vartheta_1, vartheta_2 a radius rA (entrambi
						% al medesimo centro B, stessa ampiezza, antiorario)
						\pgfmathsetmacro{\dvartheta}{atan(\ss/(2*\R))}
						
						% vartheta_1: BA → BA' (rotazione del body 1)
						\draw[->, gray!85, thin]
						({\R+\rA*cos(180)}, {\rA*sin(180)})
						arc[start angle=180, end angle={180+\dvartheta}, radius=\rA];
						\node[font=\scriptsize, gray!85, anchor=west]
						at ({\R+\rA*cos(180+\dvartheta)+0.05}, {\rA*sin(180+\dvartheta)})
						{$\vartheta_1$};
						
						% vartheta_2: BD → BD' (rotazione del body 2)
						\draw[->, gray!85, thin]
						({\R+\rA*cos(135)}, {\rA*sin(135)})
						arc[start angle=135, end angle={135+\dvartheta}, radius=\rA];
						\node[font=\scriptsize, gray!85, anchor=east]
						at ({\R+\rA*cos(135+\dvartheta)-0.05}, {\rA*sin(135+\dvartheta)})
						{$\vartheta_2$};
						
						% marker su D' (cerniera interna nella configurazione variata)
						\fill[gray!70] ({-\ss/2}, {\R-\ss/2}) circle (1.8pt);
						
						% --- etichetta pannello, sotto tutto ---
						\node[font=\scriptsize] at ({(\xV-\R)/2}, \yH-0.55) {(b)};
						
					\end{scope}
					
				\end{tikzpicture}
			}
			\caption{Soluzione grafica per i cedimenti esterni in $A$.
				\emph{(a)}~cedimento $s_1$ (orizzontale): tre centri distinti,
				$C_1$ all'intersezione fra la verticale per $A$ e la retta $BD$.
				\emph{(b)}~cedimento $s_2$ (verticale): caso degenere,
				$C_1\equiv C_2\equiv B$, e il sistema ruota rigidamente
				attorno a $B$.}
			\label{fig:arco_cedimenti_esterni}
		\end{figure}
